\documentclass[aps,prl,reprint,amsmath,amssymb,nofootinbib,raggedbottom]{revtex4-2}
\usepackage[T1]{fontenc}
\usepackage{tnnotes,booktabs}
\usepackage[colorlinks=true,citecolor=blue,linkcolor=blue,urlcolor=blue]{hyperref}
\hypersetup{pdftitle={A Diagrammatic Proof of the Haldane Gap via Spacetime Tensor Networks}}

% ------------------------------------------------ diagrams specific to this note
% All other diagrams are written inline with \tnchain, \tncolumn and prooflanes
% from tnnotes.sty.
% The local MPO tensor with its four legs named.
\newcommand{\Wlegs}{\tikz[tn]{\draw (-.95,0)--(.95,0); \draw (0,-.95)--(0,.95);
 \node[mpo] at (0,0) {$W$};
 \node[anchor=east] at (-1,0) {$a$}; \node[anchor=west] at (1,0) {$b$};
 \node[anchor=south] at (0,1) {$s$}; \node[anchor=north] at (0,-1) {$s'$};}}
% Folded contraction: the upper row is a ket, the lower row its adjoint.
\newcommand{\folded}[1]{\tikz[tn,tn centered]{
 \begin{scope}[yshift=.8em]\tnchainpath[capr=v_j,name=u]{#1}\end{scope}
 \begin{scope}[yshift=-.8em]\tnchainpath[capr=\bar v_j,name=d]{#1}\end{scope}
 \draw[rounded corners=3pt] (u-L)--++(-.35em,0)|-(d-L);}}
% Two complete thermal MPOs stacked on their physical legs. The equality
% (A13) is of full operators; no sitewise fusion is asserted.
\newcommand{\stackedmpo}{\tikz[tn,tn centered]{
 \begin{scope}[yshift=.95em]
  \tnchainpath[node=mpo,legs=.95em,caps,name=u]{W_\beta,W_\beta,\cdots,W_\beta}\end{scope}
 \begin{scope}[yshift=-.95em]
  \tnchainpath[node=mpo,legs=.95em,caps,name=d]{W_\beta,W_\beta,\cdots,W_\beta}\end{scope}}}

\begin{document}
\title{A Diagrammatic Proof of the Haldane Gap via Spacetime Tensor Networks}
\date{October 8, 2026}
\begin{abstract}
Following the recent machine-assisted resolution of the Haldane conjecture for periodic and boundary-field spin-1 Heisenberg chains by OpenAI~\cite{source,boundary}, we show that the analytical derivations can be recast as elementary tensor-network contractions.
Formulating the system as an exact continuous-time thermal matrix product operator exposes the underlying spacetime geometry.
The step from a finite-scale estimate to a uniform gap then reduces to a contractive scale-doubling recurrence on planar diagrams, under three pages long and formalized in Lean~4 with the two numerical inputs as its only unproved hypotheses.
The finite-scale estimate itself is taken unchanged from Ref.~\cite{source}, whose certificate uses spectral filters and character projections.
This illustrates how diagrammatic representations can clarify and compress machine-generated proofs in quantum many-body physics.
\end{abstract}
\maketitle
\setlength{\abovedisplayskip}{5pt plus 2pt minus 1pt}
\setlength{\belowdisplayskip}{5pt plus 2pt minus 1pt}

We consider the antiferromagnetic spin-1 Heisenberg chain with periodic
boundary conditions $\bm S_{L+1}=\bm S_1$ and $J=k_B=1$:
\begin{equation}
 H_L=\sum_{j=1}^L\bm S_j\cdot\bm S_{j+1},\qquad
 Z_L(\beta)=\Tr e^{-\beta H_L}.
 \label{eq:model}
\end{equation}
We take even $L\ge4$ so the ring is bipartite and the transfer trace is an
even power. Let $E_0(L)\le E_1(L)\le\cdots$ include multiplicities, with gap
$\Delta_L=E_1(L)-E_0(L)$ (\levels) and Gibbs weights
$w_m=e^{-\beta E_m(L)}/Z_L(\beta)$:
\begin{equation}
 \frac{w_1}{w_0}=e^{-\beta\Delta_L}\le e^{-c\beta}
 \quad\Longrightarrow\quad \Delta_L\ge c.
 \label{eq:target}
\end{equation}
The ratio decays in $\beta$ like \decayglyph; the task is to make its
rate $c$ independent of $L$:
\begin{theorem}[Periodic gap~\cite{source}]\label{thm:gap}
For every even $L\ge2304$, the chain~\eqref{eq:model} has a unique ground
state and $\Delta_L>(\log20)/784\simeq0.003821$.
\end{theorem}

Theorem~\ref{thm:gap} was established by OpenAI~\cite{source} via a thirty-page coordinate construction using degree-12 spectral filters and dihedral character decompositions. After showing how a thermal purity near one forces a gap, our proof has three stages: (i) an exact continuous-time thermal MPO yields a self-adjoint spatial transfer operator; (ii) spacetime contraction identities couple spatial and thermal purities so joint scale doubling quadratically suppresses both defects (the \emph{spacetime bootstrap}); and (iii) interpolating between successive doubled lengths covers all even $L\ge2304$.

Why tensor-network diagrams? In Eq.~\eqref{eq:model}, $L$ counts sites in a sum whereas $\beta$ sits in an exponential. Drawing $Z_L(\beta)$ as a two-dimensional tensor network~\cite{tang,pirvu} puts space and imaginary time on equal footing on a torus: vertical double legs \physleg{} carry the physical spin space through imaginary time (trace $\Tr$; blue, thermal \tikz[tn]{\node[gibbs]{$\rho$};}), horizontal legs \virtleg{} couple neighboring sites in the virtual bond space (trace $\trv$; red, spatial \tikz[tn]{\node[xop]{$\sigma$};}), joined legs \tnchain[node=xvertex]{,} are contracted, and closed loops \tnchain[node=xvertex,ring]{{}} are traces. Reading the same network along either axis turns the proof into planar diagram identities; Appendix~\ref{app:open} caps the horizontal ends with boundary vectors for an open chain.

\emph{From purity to gap.---}
For a positive trace-one operator with eigenvalues $p_0\ge p_1\ge\cdots$, the
purity obeys
\begin{equation}
 P=\sum_i p_i^2\le p_0\sum_i p_i=p_0.
 \label{eq:weight}
\end{equation}
Thus the \emph{purity defect} $1-P$ bounds $1-p_0\le1-P$: a purity near one
forces a spectrum like \specbars{.86,.08,.04,.02} with a single dominant bar.
For the Gibbs state, $P>1/2$ implies a unique ground state, with
$1-w_0\le1-P$ and $e^{-\beta\Delta_L}=w_1/w_0\le(1-P)/P$. To reach
Eq.~\eqref{eq:target} it suffices to show that $1-P$ falls like $e^{-c\beta}$
uniformly in $L$. Squaring in imaginary time improves purity at fixed $L$
[Eq.~\eqref{eq:squaring}], but growing $L$ adds thermal excitations; advancing
$L$ and $\beta$ together requires a spatial transfer operator.

\emph{Thermal MPO.---}
For the classical Ising chain, $Z_n=\Tr K^n$ with a $2\times2$ transfer matrix
$K$ stepping along the ring; here we construct a real symmetric operator
$X_\beta$ on the virtual bond space with $Z_n(\beta)=\trv X_\beta^n$. In the Cartesian basis
$s,s'\in\{x,y,z\}$, $(S^\alpha)_{ss'}=-i\epsilon_{\alpha ss'}$ gives real
operators $G_\alpha=iS^\alpha$, $\|G_\alpha\|=1$, and
$-H_L=\sum_{j,\alpha}G_{\alpha,j}G_{\alpha,j+1}$. Since the simplex
$0<\tau_1<\cdots<\tau_k<\beta$ in imaginary time $\tau$ has volume
$\beta^k/k!$, $e^{\beta A}=\sum_{k\ge0}\int_{0<\tau_1<\cdots<\tau_k<\beta}d\tau_1\cdots d\tau_k\,A^k$.
For $A=-H_L$, each factor selects a bond and label $\alpha$, inserting
$G_\alpha$ on both sites of that bond at time $\tau$. For bond events
$(\tau,\alpha)$ in $a,b$ on the left and right bonds of a site, the local MPO
tensor multiplies the matrices hitting its worldline in time order ($\mathcal T_\tau$ placing later times to the left):
\begin{equation}
 \Wlegs=W_\beta^{ss'}(a,b)=
 \left\langle s\left|\mathcal T_\tau
 \prod_{(\tau,\alpha)\in a\cup b}G_\alpha\right|s'\right\rangle.
 \label{eq:local}
\end{equation}
For $a=\{(\tau_1,x),(\tau_3,y)\}$ and $b=\{(\tau_2,z)\}$ with
$\tau_1<\tau_2<\tau_3$, the site worldline and its physical trace
\Wtrace${}={}$\Xvertex{} are
\begin{proofdiagram}
\begin{tikzpicture}[x=1cm,y=.62cm,font=\scriptsize,line width=.45pt]
 % Imaginary-time axis on the far left
 \draw[-{Latex[length=1.2mm]},mpoblue] (-.35,.2)--(-.35,1.75)
   node[above] {$\tau$};
 \draw[mpoblue] (-.41,.2)--(-.29,.2) node[left=2pt,text=black!65] {$0$};
 \draw[mpoblue] (-.41,1.75)--(-.29,1.75) node[left=2pt,text=black!65] {$\beta$};
 % Open worldline inside W_beta^{ss'}(a,b)
 \fill[mpoblue!12,rounded corners=2pt] (.7,.25) rectangle (1.9,1.7);
 \draw[mpoblue!70,densely dashed,rounded corners=2pt] (.7,.25) rectangle (1.9,1.7);
 \node[mpoblue!80!black,anchor=south west,inner sep=1pt] at (.7,1.73) {$W$};
 \draw[line width=.7pt] (1.3,.05)--(1.3,1.9);
 \node[anchor=south] at (1.3,1.9) {$\langle s|$};
 \node[anchor=north] at (1.3,.05) {$|s'\rangle$};
 % Rungs from left bond a (tau_1, tau_3) and right bond b (tau_2)
 \draw[mpored,line width=.65pt] (.15,.55)--(1.3,.55);
 \fill[mpored] (1.3,.55) circle[radius=1.3pt];
 \node[anchor=east,text=black!75] at (.13,.55) {$\tau_1$};
 \node[anchor=south west,inner sep=1pt] at (1.34,.55) {$G_x$};
 \draw[mpoblue!85!black,line width=.65pt] (1.3,.98)--(2.45,.98);
 \fill[mpoblue!85!black] (1.3,.98) circle[radius=1.3pt];
 \node[anchor=west,text=black!75] at (2.47,.98) {$\tau_2$};
 \node[anchor=south east,inner sep=1pt] at (1.26,.98) {$G_z$};
 \draw[mpored,line width=.65pt] (.15,1.4)--(1.3,1.4);
 \fill[mpored] (1.3,1.4) circle[radius=1.3pt];
 \node[anchor=east,text=black!75] at (.13,1.4) {$\tau_3$};
 \node[anchor=south west,inner sep=1pt] at (1.34,1.4) {$G_y$};
 \node[mpored!85!black,anchor=south] at (.42,1.45) {$a$};
 \node[mpoblue!85!black,anchor=south] at (2.18,1.03) {$b$};
 \node[anchor=north,font=\footnotesize] at (1.3,-.18)
   {$W_\beta^{ss'}(a,b)=\langle s|G_yG_zG_x|s'\rangle$};
 % Traced worldline X_beta(a,b) on the right
 \fill[mpored!10,rounded corners=2pt] (4.55,.25) rectangle (5.75,1.7);
 \draw[mpored!70,densely dashed,rounded corners=2pt] (4.55,.25) rectangle (5.75,1.7);
 \node[mpored!80!black,anchor=south west,inner sep=1pt] at (4.55,1.73) {$X_\beta$};
 \draw[line width=.7pt] (5.15,.25)--(5.15,1.7);
 % Trace loop closing top to bottom
 \draw[line width=.55pt,black!65] (5.15,1.7)--(5.15,1.82)
   arc[start angle=180,end angle=0,radius=.38]
   --(5.91,.13) arc[start angle=0,end angle=-180,radius=.38]--(5.15,.25);
 \draw[mpored,line width=.65pt] (4.0,.55)--(5.15,.55);
 \fill[mpored] (5.15,.55) circle[radius=1.3pt];
 \node[anchor=south west,inner sep=1pt] at (5.19,.55) {$G_x$};
 \draw[mpoblue!85!black,line width=.65pt] (5.15,.98)--(6.3,.98);
 \fill[mpoblue!85!black] (5.15,.98) circle[radius=1.3pt];
 \node[anchor=south east,inner sep=1pt] at (5.11,.98) {$G_z$};
 \draw[mpored,line width=.65pt] (4.0,1.4)--(5.15,1.4);
 \fill[mpored] (5.15,1.4) circle[radius=1.3pt];
 \node[anchor=south west,inner sep=1pt] at (5.19,1.4) {$G_y$};
 \node[mpored!85!black,anchor=south] at (4.27,1.45) {$a$};
 \node[mpoblue!85!black,anchor=south] at (6.15,1.03) {$b$};
 \node[anchor=north,font=\footnotesize] at (5.2,-.18)
   {$X_\beta(a,b)=\Tr(G_yG_zG_x)$};
 % Arrow between them
 \draw[flow,black!60] (2.95,.98)--(3.65,.98)
   node[midway,above,text=black!70] {$\sum_s$};
\end{tikzpicture}
\end{proofdiagram}
Operators on different sites commute, so once the events on every bond are
fixed, the time-ordered product over the ring splits into one worldline per
site: the bond history plays the role of the Ising spin between two copies of
$K$. Summing the three labels $\alpha\in\{x,y,z\}$ and integrating ordered
times defines a measure $\mu_\beta$ of mass $\sum_{k\ge0}(3\beta)^k/k!=e^{3\beta}$
(unit mass on the empty history).
Contracting $n$ tensors around the ring recovers the convergent exponential
series as an exact MPO on $L^2(\mu_\beta)$ without bond truncation~\cite{tang,pirvu}:
\begin{equation}
 e^{-\beta H_n}=\tnchain[node=mpo,legs,ring]{W,W,\cdots,W,W}.
 \label{eq:mpo}
\end{equation}
Because $a$ and $b$ enter Eq.~\eqref{eq:local} through $a\cup b$ and each event
inserts the same real matrix $G_\alpha$ at both ends of its bond (\emph{bond
reflection positivity}, $-\bm S_j\cdot\bm S_{j+1}=\sum_\alpha G_{\alpha,j}G_{\alpha,j+1}$),
the spatial transfer kernel $X_\beta(a,b)=\sum_s W_\beta^{ss}(a,b)$ is real and
symmetric ($X_\beta(a,b)=X_\beta(b,a)$). Since $\|G_\alpha\|=1$ on $\mathbb C^3$,
$|X_\beta(a,b)|\le\Tr I=3$; with the finite mass $e^{6\beta}$ of
$\mu_\beta\times\mu_\beta$, $X_\beta$ is self-adjoint and Hilbert--Schmidt on
$L^2(\mu_\beta)$, giving a discrete real spectrum
$|\lambda_0(\beta)|\ge|\lambda_1(\beta)|\ge\cdots$ ($\sum_i\lambda_i(\beta)^2<\infty$).
Closing the ring of transfer vertices gives, for even $n$,
\begin{equation}
 Z_n(\beta)=\tnchain[node=xvertex,ring]{,,\cdots,,}=\trv X_\beta^n=\sum_i|\lambda_i(\beta)|^n,
 \label{eq:moment}
\end{equation}
so raising the eigenvalues to the $n$th power concentrates the weight on
$\lambda_0$:
\specbars[mpored]{.5,.3,.2}${}\mapsto{}$\specbars[mpored]{.87,.11,.02} ($n=4$).

\emph{Spatial and thermal purities.---}
The transfer operator gives space its own density operator. Cut the ring of
Eq.~\eqref{eq:moment} at one bond,
\tnchain[node=xvertex,ring]{,\cdots,}${}\mapsto{}$\tnchain[node=xvertex]{,\cdots,},
and normalize both the transfer and Gibbs operators:
\begin{equation}
 \sigma_n(\beta)=\frac{\tnchain[node=xvertex]{,,\cdots,,}}{Z_n(\beta)}=\frac{X_\beta^n}{Z_n(\beta)},\qquad
 \rho_n(\beta)=\frac{e^{-\beta H_n}}{Z_n(\beta)}.
 \label{eq:states}
\end{equation}
For even $n$, both operators are positive with unit trace; denote their
eigenvalues by $p_i=|\lambda_i(\beta)|^n/Z_n(\beta)$ and $w_m$, respectively.
Their purities join two copies, along
space for $\sigma_n$ and along imaginary time for $\rho_n$. Drawing
$Z_n(\beta)$ as a spacetime torus \Zbox[.7]21{} of width $n$ and height $\beta$
(dashed edges identified), both are one ratio of tori:
\begin{equation}
 \begin{aligned}
 \Ps(n,\beta)&=\trv\sigma_n^2=\tnchain[ring]{\sigma_n,\sigma_n}
   =\frac{Z_{2n}(\beta)}{Z_n(\beta)^2}=\frac{\Zt41}{\Ztsq21},\\[2pt]
 \Pt(n,\beta)&=\Tr\rho_n^2\;=\tncolumn[ring]{\rho_n,\rho_n}
   =\frac{Z_n(2\beta)}{Z_n(\beta)^2}=\frac{\Zt22}{\Ztsq21}.
 \end{aligned}
 \label{eq:purities}
\end{equation}

\emph{Squaring.---}
Squaring improves a purity: squaring $\sigma_n$ doubles the length; squaring
$\rho_n$ doubles the inverse temperature:
\begin{equation}
 \tnchain{\sigma_n,\sigma_n}=\Ps(n,\beta)\,\tnchain{\sigma_{2n}},\quad
 \tncolumn{\rho_n,\rho_n}=\Pt(n,\beta)\,\tncolumn[item={minimum height=1.6em}]{\rho_n(2\beta)}.
 \label{eq:squaring}
\end{equation}
Both replace the normalized eigenvalues $p_i$ by $p_i^2/P$,
concentrating the spectrum: \specbars{.5,.3,.2}${}\mapsto{}$\specbars{.66,.24,.1}.
Since $1-P=2\sum_{i<k}p_ip_k$, dropping cross terms from
$(\sum_{i<k}p_ip_k)^2$ bounds the new defect $1-P'$ via
$f(z):=z^2/[2(1-z)^2]$:
\begin{equation}
 1-P'=\frac{2\sum_{i<k}p_i^2p_k^2}{P^2}
 \le\frac{(1-P)^2}{2P^2}=f(1-P).
 \label{eq:square}
\end{equation}
Because $f$ increases on $[0,1)$ with $f(z)\sim z^2/2$, squaring
roughly squares the purity defect $1-P$, keeping it far below the
diagonal, \fcurve{} for $0\le z\le.3$.

\emph{Spacetime coupling identities.---}
Each squaring doubles one side of the torus: space ($n\to 2n$), \Zt21${}\mapsto{}$\Zt41, or imaginary time ($\beta\to 2\beta$), \Zt21${}\mapsto{}$\Zt22. Because a uniform gap needs both sides to grow together, we need identities trading purity in one direction for purity in the other.
Contracting four copies of $Z_n(\beta)$ in either spatial or temporal order yields the composite torus $Z_{2n}(2\beta)$ (red arrows join copies in space, blue in time, labeled by the purity factor they multiply by):
\begin{proofdiagram}
\begin{tikzpicture}[x=1cm,y=1cm,font=\footnotesize]
 % One n x beta torus = 2 x 1 cells; gaps close as copies are joined.
 \def\c{.3}\def\g{.14}
 \foreach \a in {0,1}{\foreach \b in {0,1}{                       % Z_n(beta)^4
   \tnpatch{{.6+\a*(2*\c+\g)}}{{2.28+\b*(\c+\g)}}{2}{1}{\c}}}
 \foreach \b in {0,1}{\tnpatch{5.4}{{2.28+\b*(\c+\g)}}{4}{1}{\c}}  % Z_2n(beta)^2
 \foreach \a in {0,1}{\tnpatch{{.6+\a*(2*\c+\g)}}{.5}{2}{2}{\c}}  % Z_n(2beta)^2
 \tnpatch{5.4}{.5}{4}{2}{\c}                                     % Z_2n(2beta)
 \node[anchor=north] at (1.27,2.23) {$Z_n(\beta)^4$};
 \node[anchor=north] at (6.0,2.23) {$Z_{2n}(\beta)^2$};
 \node[anchor=north] at (1.27,.45) {$Z_n(2\beta)^2$};
 \node[anchor=north] at (6.0,.45) {$Z_{2n}(2\beta)$};
 \draw[flow,mpored] (2.15,2.65)--(5.2,2.65)
   node[midway,above,text=black] {$\times\Ps(n,\beta)^2$};
 \draw[flow,mpored] (2.15,.80)--(5.2,.80)
   node[midway,above,text=black] {$\times\Ps(n,2\beta)$};
 \draw[flow,mpoblue] (1.27,1.82)--(1.27,1.26)
   node[midway,right,text=black] {$\times\Pt(n,\beta)^2$};
 \draw[flow,mpoblue] (6.0,1.82)--(6.0,1.26)
   node[midway,left,text=black] {$\times\Pt(2n,\beta)$};
 \draw[flow,mpored] (-.45,.5)--(.05,.5) node[right,text=black] {$L$};
 \draw[flow,mpoblue] (-.45,.5)--(-.45,1.0) node[above,text=black] {$\beta$};
\end{tikzpicture}
\end{proofdiagram}
Writing each purity as its ratio of tori, Eq.~\eqref{eq:purities}, the
two paths cancel to the same quotient:
\begin{equation}
 \underbrace{\frac{\Zt42}{\Ztsq22}}_{\Ps(n,2\beta)}
 \underbrace{\frac{\Ztsq22}{\Ztsq[4]21}}_{\Pt(n,\beta)^2}
 =\frac{\Zt42}{\Ztsq[4]21}=
 \underbrace{\frac{\Ztsq41}{\Ztsq[4]21}}_{\Ps(n,\beta)^2}
 \underbrace{\frac{\Zt42}{\Ztsq41}}_{\Pt(2n,\beta)}.
 \label{eq:cancel}
\end{equation}
Solving for one factor, and repeating at length $2n$, gives
\begin{align}
 \Ps(n,2\beta)&=\frac{\Ps(n,\beta)^2\Pt(2n,\beta)}{\Pt(n,\beta)^2},
 \label{eq:couple1}\\
 \Pt(4n,\beta)&=\frac{\Pt(2n,\beta)^2\Ps(2n,2\beta)}{\Ps(2n,\beta)^2}.
 \label{eq:couple2}
\end{align}

\emph{Doubling estimate.---}
Alternating the two coupling identities with the two squarings doubles the
torus in both directions (\torusseq) and roughly squares the purity defect. For
even $n$ and $0\le\varepsilon\le1/100$,
\begin{equation}
 \begin{gathered}
 \Ps(n,\beta),\ \Pt(2n,\beta)\ge1-\varepsilon\\
 \Longrightarrow\quad
 \Ps(2n,2\beta),\ \Pt(4n,2\beta)\ge1-5\varepsilon^2.
 \end{gathered}
 \label{eq:update}
\end{equation}
Four moves prove Eq.~\eqref{eq:update}. Read the diagram below from top to
bottom (red boxes bound $\Ps$, blue $\Pt$; solid given or proved, dashed
intermediate; circles label the identity and grey arrows supply the second
input):
\begin{proofdiagram}
\begin{prooflanes}\def\lanedy{.78}
 \lanebound{s}{s0}{0}{0}{$\Ps(n,\beta)\ge1-\varepsilon$}
 \lanebound{t}{t0}{1}{0}{$\Pt(2n,\beta)\ge1-\varepsilon$}
 \lanebound{so}{s1}{0}{1}{$\Ps(n,2\beta)\ge1-3\varepsilon$}
 \lanebound{s}{s2}{0}{2}{$\Ps(2n,2\beta)\ge1-5\varepsilon^2$}
 \lanebound{to}{t1}{1}{3}{$\Pt(4n,\beta)\ge1-3\varepsilon$}
 \lanebound{t}{t2}{1}{4}{$\Pt(4n,2\beta)\ge1-5\varepsilon^2$}
 \lanemove{s0}{s1}{1}{\eqref{eq:couple1}} \laneuse{t0}{1}
 \lanemove[mpored]{s1}{s2}{2}{square}
 \lanemove{t0}{t1}{3}{\eqref{eq:couple2}} \laneuse{s2}{3}
 \lanemove[mpoblue]{t1}{t2}{4}{square}
 \lanetag[above]{s0}{given} \lanetag[above]{t0}{given}
 \lanetag[below]{s2}{proved} \lanetag{t2}{proved}
\end{prooflanes}
\end{proofdiagram}
The proved pair is the given pair at $(2n,2\beta)$, so Eq.~\eqref{eq:update} iterates.

\emph{Move 1} (coupling). Since $\Pt(n,\beta)\le1$,
Eq.~\eqref{eq:couple1} gives $\Ps(n,2\beta)\ge(1-\varepsilon)^3\ge1-3\varepsilon$.
\emph{Move 2} (squaring in space). For even $n$, $\sigma_n(2\beta)$ has
nonnegative spectrum and squares to $\sigma_{2n}(2\beta)$. Since
$1-(1-\varepsilon)^3\le3\varepsilon$ and $9/[2(1-3\varepsilon)^2]\le5$ for
$\varepsilon\le1/100$, Eq.~\eqref{eq:square} and monotonicity of $f$ give
$\Ps(2n,2\beta)\ge1-\delta$, with
\begin{equation}
 \delta=f\bigl(1-(1-\varepsilon)^3\bigr)
 =\tfrac12\bigl[(1-\varepsilon)^{-3}-1\bigr]^2
 \le5\varepsilon^2\le\varepsilon.
 \label{eq:delta}
\end{equation}
\emph{Move 3} (coupling). Equation~\eqref{eq:couple2} with
$\Ps(2n,\beta)\le1$ and $\delta\le\varepsilon$ gives
$\Pt(4n,\beta)\ge(1-\varepsilon)^2(1-\delta)\ge(1-\varepsilon)^3$.
\emph{Move 4} (squaring in time). Squaring $\rho_{4n}(\beta)$ to
$\rho_{4n}(2\beta)$ gives $\Pt(4n,2\beta)\ge1-\delta\ge1-5\varepsilon^2$,
proving Eq.~\eqref{eq:update}.

\emph{Starting estimate and iteration.---}
The doubling estimate needs one starting scale where both defects are below
$1/100$. Proposition~5.5 of Ref.~\cite{source} gives the $j=0$ starting estimates:
\begin{equation}
 \Ps(2304,784),\ \Pt(4608,784)\ge0.99.
 \label{eq:input}
\end{equation}
In Ref.~\cite{source}, shifting $H_L\to H_L+\kappa L I$ ($\kappa=700741/500000$)
normalizes $\lambda_0>1$; the factor $e^{-\kappa\beta L}$ in $Z_L(\beta)$
cancels from both purity ratios~\eqref{eq:purities}.
After $j\ge0$ doublings, set
\begin{equation}
 n_j=2304\,2^j,\qquad \beta_j=784\,2^j,\qquad
 \varepsilon_j=\frac{20^{-2^j}}5.
 \label{eq:scales}
\end{equation}
Since $\varepsilon_0=1/100$ and $\varepsilon_{j+1}=5\varepsilon_j^2$,
Eq.~\eqref{eq:update} gives
$\Ps(n_j,\beta_j),\Pt(2n_j,\beta_j)\ge1-\varepsilon_j$ by induction on $j$:
each inductive step runs Moves~1--4 once.
The defect falls doubly exponentially in $j$,
$\varepsilon_0=10^{-2}$\,\sparkline{-22}{-2}{-2,-3.30,-5.90,-11.11,-21.52}\,$\varepsilon_4\approx3\times10^{-22}$ (log scale),
and so exponentially in $\beta_j$,
$\varepsilon_j=\frac15\exp[-\beta_j(\log20)/784]$, as needed in Eq.~\eqref{eq:target}.

\emph{All even lengths.---}
These estimates hold at the doubled lengths $n_j$ and $2n_j$. To cover every
even $L\ge2304$, choose $j$ with $n_j\le L\le2n_j$. Since the weights
$p_i=|\lambda_i(\beta)|^n/Z_n(\beta)\in[0,1]$ sum to one, any power $r\ge1$
gives $Z_{rn}(\beta)/Z_n(\beta)^r=\sum_ip_i^r\le\sum_ip_i=1$; taking
$r=L/n_j\ge1$ at $\beta_j$ and $r=2n_j/L\ge1$ at $2\beta_j$ yields
\begin{equation}
 Z_L(\beta_j)\le Z_{n_j}(\beta_j)^{L/n_j},\qquad
 Z_L(2\beta_j)\ge Z_{2n_j}(2\beta_j)^{L/(2n_j)}.
 \label{eq:moments}
\end{equation}
Dividing $Z_L(2\beta_j)$ by $Z_L(\beta_j)^2$, substituting the torus identity
$Z_{2n_j}(2\beta_j)/Z_{n_j}(\beta_j)^4 = \Pt(2n_j,\beta_j)\Ps(n_j,\beta_j)^2$
[Eq.~\eqref{eq:cancel}], and using $L/(2n_j)\le1$ with
$(1-\varepsilon_j)^3\ge1-3\varepsilon_j>\tfrac12$ gives
\begin{equation}
 \Pt(L,\beta_j)
 \ge\bigl[\Pt(2n_j,\beta_j)\Ps(n_j,\beta_j)^2\bigr]^{L/(2n_j)}
 \ge1-3\varepsilon_j.
 \label{eq:interpolate}
\end{equation}

\emph{Recovering the gap.---}
With $w_0\ge w_1\ge\cdots$ for $\rho_L(\beta_j)$, Eqs.~\eqref{eq:weight} and
\eqref{eq:interpolate} give $w_0\ge1-3\varepsilon_j>1/2$: the Gibbs spectrum
looks like \specbars{.97,.02,.01} rather than \specbars{.49,.49,.02}, so the
ground state is unique and $w_1\le 1-w_0\le 3\varepsilon_j$. Since
$3\varepsilon_j/(1-3\varepsilon_j)\le(60/97)20^{-2^j}<20^{-2^j}$,
\begin{equation}
 e^{-\beta_j\Delta_L}=\frac{w_1}{w_0}
 \le\frac{1-w_0}{w_0}
 \le\frac{3\varepsilon_j}{1-3\varepsilon_j}
 <20^{-2^j}.
 \label{eq:gapratio}
\end{equation}
Taking logarithms and dividing by $\beta_j=784\,2^j$ gives, for every
even $L\ge2304$,
\begin{equation}
 \boxed{\Delta_L>\frac{\log20}{784}\simeq0.003821.}
 \label{eq:gap}
\end{equation}
This proves Theorem~\ref{thm:gap}, conditional on Eq.~\eqref{eq:input};
in particular $\liminf_{L\to\infty,\;L\,\mathrm{even}}\Delta_L\ge(\log20)/784>0$.
The scale-doubling recurrence and the gap deduction are formally verified in Lean~4 (\path{lean/CertifiedGaps/PurityCriterion.lean}).


\section*{Discussion: Representation in Machine-Assisted Proofs}
The uniform spectral gap theorems for the periodic and boundary-field spin-1 Heisenberg chains, together with the certified starting estimates, come from the machine-assisted proofs of OpenAI~\cite{source,boundary}. Those thirty-page proofs locate eigenvalues and connect scales through explicit coordinate constructions, such as degree-12 polynomial spectral filters and dihedral character projectors, alongside lemmas unused in the final inequality: working by local symbolic rewriting in a fixed basis without the geometric pictures of many-body physics, long chains of deductions naturally build up intermediate algebra to stay sound.

Here, the authors asked GPT-6 to re-examine the argument in terms of spacetime tensor networks and worked through the resulting structure interactively. Treating imaginary time and space alike on a torus reduces the self-adjoint spatial transfer operator and the contractive scale-doubling iteration to short identities between planar diagrams. Tensor-network diagrams are thus useful not only for calculation, but also for reading, simplifying, and extracting the physical mechanism of a machine-generated proof. The authors checked all mathematical and formal proofs.

\begin{thebibliography}{4}
\bibitem{source}
OpenAI, \emph{The periodic spin-one Haldane gap}, manuscript dated
September 24, 2026, Proposition 5.5 and Appendices A--C.
\bibitem{boundary}
OpenAI, \href{https://github.com/openai/math/blob/main/preprints/A-boundary-field-gap-for-the-spin-one-Heisenberg-chain-September-24-2026/paper.pdf}
{\emph{A boundary-field gap for the spin-one Heisenberg chain}},
manuscript dated September 24, 2026, Secs.~2--6 and Appendices A and B.
\bibitem{tang}
W. Tang, H.-H. Tu, and L. Wang,
\href{https://doi.org/10.1103/PhysRevLett.125.170604}
{Continuous matrix product operator approach to finite temperature quantum states},
Phys. Rev. Lett. \textbf{125}, 170604 (2020).
\bibitem{pirvu}
B. Pirvu, V. Murg, J. I. Cirac, and F. Verstraete,
\href{https://doi.org/10.1088/1367-2630/12/2/025012}
{Matrix product operator representations},
New J. Phys. \textbf{12}, 025012 (2010).
\end{thebibliography}

\clearpage
\appendix
\setcounter{secnumdepth}{2}
\section{Open chain with boundary fields}\label{app:open}
The networks above also bound the gap of the open chain with equal
endpoint fields. On $n$ spin-1 sites, with
$\sum_\alpha(S^\alpha)^2=2I$, let
\begin{equation}
 H_n^{\same}=\sum_{i=1}^{n-1}\bm S_i\cdot\bm S_{i+1}
       -h(S_1^z+S_n^z),\qquad h=\frac35.
 \label{eq:b-model}
\end{equation}
\begin{theorem}[Boundary-field gap~\cite{boundary}]\label{thm:b-gap}
For every odd $n\ge1921$, the Hamiltonian~\eqref{eq:b-model} has a unique
ground state and $\Delta_n>(\log10)/392\simeq0.005874$.
\end{theorem}
Boundary fields are essential here: without them, each end of the open spin-1
chain carries an effective spin-$1/2$ boundary mode, leaving four nearly
degenerate low-lying states; equal fields lift this boundary degeneracy on odd
chains~\cite{boundary}. As on the ring, by Eqs.~\eqref{eq:target} and
\eqref{eq:weight}, to prove Theorem~\ref{thm:b-gap} it suffices to bound the
equal-field thermal purity $\Pt^{\same}(n,\beta)=\sum_mw_m^2$ above $1/2$:
\begin{equation}
 e^{-\beta\Delta_n}=\frac{w_1}{w_0}
 \le\frac{1-\Pt^{\same}(n,\beta)}{\Pt^{\same}(n,\beta)}.
 \label{eq:b-target}
\end{equation}

\emph{Strategy.---}
The ring proof cannot be applied directly because an open chain has no ring to
close: its partition function is not a trace of $X_\beta^n$, so there is no
spatial purity to square. Its bulk, however, shares the periodic transfer
operator $X_\beta$, and on an open chain the outer virtual legs of the first
and last sites are closed by boundary vectors (the quantum counterpart of
$u^{\mathsf T}K^{n-1}u$ for an open classical Ising chain in an endpoint field).
Whereas a periodic trace \tnchain[node=xvertex,ring]{,\cdots,} weights every
eigenvector of $X_\beta$ by one, an open chain
\tnchain[node=xvertex,caps]{,\cdots,} weights the leading eigenvalue
$\lambda(\beta)$ by the squared boundary overlap $c(\beta)$
[Eq.~\eqref{eq:b-overlap}] and vanishes if the two are orthogonal. For long
chains the open partition function is approximately $c(\beta)\lambda(\beta)^n$,
so its thermal purity is approximately
$[\lambda(2\beta)/\lambda(\beta)^2]^n c(2\beta)/c(\beta)^2$: a bulk factor
controlled by the ring times an $n$-independent boundary factor that cancels
when comparing lengths $M$ and $2M$. Doubling the length therefore costs only
the bulk factor, while squaring in imaginary time~\eqref{eq:squaring} needs no
ring. Because this requires the leading term to dominate the remainder, the
proof tracks the \emph{boundary overlap} alongside the thermal purity through a
boundary defect $e$ that shrinks as the chain grows.

\emph{Boundary vectors.---}
The open chain uses the tensor $W_\beta$ of Eq.~\eqref{eq:local}, with the
per-site energy shift $\kappa=700741/500000$ of Ref.~\cite{boundary} (called
$a$ there) absorbed into each tensor as a factor $e^{-\kappa\beta}$, i.e.,
$H_n^\nu\to H_n^\nu+\kappa nI$; this normalizes the leading periodic
eigenvalue $\lambda_0(\beta_j)>1$ at the certified scales while
canceling from every Gibbs state and purity ratio.
On an end site, the field $\pm hS^z=\mp(ih)G_z$ enters the imaginary-time
expansion as half a bond: for a history $a$ of $k$ events, define the
boundary vector on $L^2(\mu_\beta)$ by
\begin{equation}
 v_\beta(a)=
 \begin{cases}(ih)^k,&\text{all }k\text{ labels are }z,\\0,&\text{otherwise},
 \end{cases}
 \qquad \|v_\beta\|^2=e^{\beta h^2},
 \label{eq:b-vector}
\end{equation}
with value one on the empty history ($k=0$) and
$\|v_\beta\|^2=\sum_{k\ge0}(\beta h^2)^k/k!$ from the $z$-event time integrals.
The conjugated left cap $\bar v_\beta$ weights each $z$ event by
$(-ih)G_z=+hS^z$, adding $-hS_1^z$ to the Hamiltonian, whereas the unbarred
right cap $v_\beta$ weights each $z$ event by $(ih)G_z=-hS^z$ and adds
$+hS_n^z$, giving the opposite-field chain
$H_n^{\opp}:=H_n^{\same}+2hS_n^z$, \fieldopp. The $\pi$ rotation
$P=\operatorname{diag}(1,-1,-1)$ about $x$ flips $S^z\to-S^z$ and every $z$
label, so its virtual action $U_P$ reverses the right field and restores equal
fields, \fieldsame.
Opening the ring of Eq.~\eqref{eq:mpo} and capping its ends therefore gives
\begin{equation}
 e^{-\beta H_n^{\same}}=\tnchain[node=mpo,legs,capl=\bar v_\beta,capr=U_Pv_\beta]{W,W,\cdots,W,W},
 \label{eq:b-mpo}
\end{equation}
and tracing the physical legs onto the transfer chain $X_\beta^n$ of
Eq.~\eqref{eq:states} gives the open-chain partition functions (the open
counterparts of $Z_n$ in Eq.~\eqref{eq:moment}):
\begin{align}
 Y_n(\beta)&=\Tr e^{-\beta H_n^{\opp}}
 =\tnchain[capl={\bar v_\beta},capr={v_\beta}]{X_\beta^n},\quad\fieldopp,\label{eq:b-y}\\
 Y_n^{\same}(\beta)&=\Tr e^{-\beta H_n^{\same}}
 =\tnchain[capl={\bar v_\beta},capr={U_Pv_\beta}]{X_\beta^n},\quad\fieldsame.
 \label{eq:b-ys}
\end{align}
The opposite-field chain has identical bra and ket caps, so its even moments
are nonnegative; this is why the proof works with $H_n^{\opp}$ and transfers to
$H_n^{\same}$ at the end. Simultaneous spin rotations of the histories act on
$L^2(\mu_\beta)$ as real orthogonal operators commuting with $X_\beta$; in
particular $[U_P,X_\beta]=0$. Before the $\kappa$ shift, the absolute series
for the traced expansion sums to at most $3^n e^{\beta[3(n-1)+2h]}<\infty$, so
Fubini's theorem justifies regrouping the time-ordered series into the
site-by-site contractions~\eqref{eq:b-y}--\eqref{eq:b-ys}.

\emph{Transfer operator spectrum.---}
The open-chain argument needs the following facts about $X_\beta$, all supplied
by the ring. At successive inverse temperatures $\beta_{j+1}=2\beta_j$, the
ring estimates of Appendix~\ref{app:b-periodic} show that the leading
eigenvalue $\lambda_j:=\lambda_0(\beta_j)$ of Eq.~\eqref{eq:moment} is simple
and positive. Normalizing $\lambda_j$ to one splits $X_{\beta_j}/\lambda_j$ into
its leading projector $\Pi_j$ and orthogonal remainder $R_j$
($\Pi_jR_j=R_j\Pi_j=0$):
\begin{equation}
 T_j:=X_{\beta_j}/\lambda_j=\Pi_j+R_j,\qquad
 \Pi_j=|\psi_j\rangle\langle\psi_j|,\quad\|\psi_j\|=1,
 \label{eq:b-split}
\end{equation}
so $T_j^m=\Pi_j+R_j^m$, with $\|R_j\|<1$, $U_P\Pi_j=\Pi_j$, and
$0<r_j:=\lambda_{j+1}/\lambda_j^2\le1$ [Eqs.~\eqref{eq:b-tail}--\eqref{eq:b-ratios}].

\emph{Boundary overlap and residual decomposition.---}
We now split the open-chain contraction into the leading part, which carries the
boundary overlap, and a remainder. Put $v_j=v_{\beta_j}$. Inserting the leading
projector disconnects the network:
\begin{equation}
 \underbrace{\tnchain[capl={\bar v_j},capr={v_j}]{\Pi_j}}_{c_j}
 =\tnchain[capl=\bar v_j,nostubr]{\psi_j}\;\;\tnchain[nostubl,capr=v_j]{\bar\psi_j}
 =|\langle\psi_j,v_j\rangle|^2.
 \label{eq:b-overlap}
\end{equation}
The overlap $c_j$ is the squared projection of $v_j$ onto $\psi_j$,
\overlapglyph; unlike a periodic trace, where $\trv\Pi_j=1$, it vanishes when
the two are orthogonal. For even $m\ge2$, $T_j^m=\Pi_j+R_j^m$ splits the
normalized partition function $B_m(j):=Y_m(\beta_j)/\lambda_j^m>0$ into $c_j$
and a remainder $D_m(j)$:
\begin{equation}
 \underbrace{\tnchain[capl={\bar v_j},capr={v_j}]{T_j^m}}_{B_m(j)}
 =\underbrace{\tnchain[capl={\bar v_j},capr={v_j}]{\Pi_j}}_{c_j}
 +\underbrace{\tnchain[capl={\bar v_j},capr={v_j}]{R_j^m}}_{D_m(j)}.
 \label{eq:b-boundary-split}
\end{equation}
Since $R_j$ is self-adjoint, $R_j^m=R_j^{m/2}R_j^{m/2}$ folds the chain at its
midpoint into a squared norm, showing that the remainder is nonnegative:
\begin{equation}
 D_m(j)=\folded{R_j^{m/2}}=\|R_j^{m/2}v_j\|^2\ge0.
 \label{eq:b-fold}
\end{equation}
The boundary defect is defined by
\begin{equation}
 e_m(j):=\frac{D_m(j)}{B_m(j)}=1-\frac{c_j}{B_m(j)},
 \label{eq:b-e}
\end{equation}
the grey fraction of the bar, \stackbar{.14}, drawn at the starting
value $e_{142}(2)\le.14$ of Eq.~\eqref{eq:b-seeds}. Any $e_m<1$ guarantees a
nonzero overlap with the leading eigenvector.

\emph{Spatial remainder suppression.---}
A longer chain shrinks the remainder while leaving the leading overlap $c_j$
unchanged. For even $m\ge k$, moving $k/2$ factors of $R_j$ onto each boundary
vector gives $u=R_j^{k/2}v_j$ and
\begin{equation}
 D_m(j)=\tnchain[capl={\bar u},capr={u}]{R_j^{m-k}}
 \le\|R_j\|^{m-k}\underbrace{\|u\|^2}_{D_k(j)},
 \label{eq:b-residual}
\end{equation}
decaying geometrically by $\|R_j\|^{m-k}$ \decayglyph[black!50] and ensuring
$c_j\le B_m(j)\le B_k(j)$.

\emph{Thermal operator composition.---}
Because Eq.~\eqref{eq:b-mpo} is an exact operator identity on the physical
Hilbert space, composing $e^{-\beta H_m^{\opp}}e^{-\beta H_m^{\opp}}=e^{-2\beta H_m^{\opp}}$
in imaginary time couples temperatures:
\begin{equation}
 \underbrace{\stackedmpo}_{e^{-\beta H_m^{\opp}}\,e^{-\beta H_m^{\opp}}}
 =\underbrace{\tnchain[node=mpo,legs=1.9em,caps,item={minimum height=1.3em}]{W_{2\beta},W_{2\beta},\cdots,W_{2\beta}}}_{e^{-2\beta H_m^{\opp}}}.
 \label{eq:b-stack-mpo}
\end{equation}
Normalizing the trace defines the opposite-field Gibbs state
$\rho_m^{\opp}(\beta)=e^{-\beta H_m^{\opp}}/Y_m(\beta)$ and its thermal purity:
\begin{equation}
 \Pt^{\opp}(m,\beta_j):=\frac{Y_m(\beta_{j+1})}{Y_m(\beta_j)^2}=\Tr\rho_m^{\opp}(\beta_j)^2\in(0,1].
 \label{eq:b-a}
\end{equation}
Substituting $Y_m(\beta_j)=\lambda_j^m B_m(j)$ and $\lambda_{j+1}=r_j\lambda_j^2$
into Eq.~\eqref{eq:b-a} links adjacent temperatures:
\begin{equation}
 B_m(j+1)=\Pt^{\opp}(m,\beta_j)r_j^{-m}B_m(j)^2.
 \label{eq:b-bridge}
\end{equation}

\emph{Coupled overlap recurrence.---}
This is Move~1: a small boundary defect at scale $j$ gives a small one at scale
$j+1$. Assuming $e_k(j)<1$ at scale $j$, let primes denote $j+1$, choose even
$m>k\ge2$, and let $\|R_{j+1}\|^{m-k}\le q<1$. Spatial suppression at $j+1$
gives $D_m'\le qD_k'$. Since $c'$ is length-independent,
$D_k'-D_m'=B_k'-B_m'$; substituting $D_k'=D_m'+B_k'-B_m'$ into $D_m'\le qD_k'$
gives the first inequality, and dividing by $B_m'$ the second:
\begin{equation}
 (1-q)D_m'\le q(B_k'-B_m'),\qquad
 e_m'\le\frac{q}{1-q}\left(\frac{B_k'}{B_m'}-1\right).
 \label{eq:b-new-e}
\end{equation}
Using $\Pt^{\opp}(k,\beta_j),r_j\le1$ and $B_m(j)\ge c_j=[1-e_k(j)]B_k(j)>0$,
composition~\eqref{eq:b-bridge} bounds the ratio at the previous temperature:
\begin{align}
 \frac{B_k'}{B_m'}
 &=\frac{\Pt^{\opp}(k,\beta_j)}{\Pt^{\opp}(m,\beta_j)}\left[\frac{B_k(j)}{B_m(j)}\right]^2r_j^{m-k}\notag\\
 &\le\frac{1}{\Pt^{\opp}(m,\beta_j)[1-e_k(j)]^2}.
 \label{eq:b-new-ratio}
\end{align}
Because only the previous overlap appears in the denominator, $e_m'<1$ establishes that the updated overlap remains positive.
When $e_k(j)$ and $1-\Pt^{\opp}(m,\beta_j)$ are $\mathcal O(x)$, choosing $m-k$ with $q\le\mathcal O(x)$ ensures that $e_m'$ contracts quadratically as $\mathcal O(x^2)$.

\emph{Thermal purity propagation.---}
This is Move~2: with the overlap under control, the length can double at almost
no cost in purity. Since $B_{2M}(j+1)\ge c_{j+1}=[1-e_M(j+1)]B_M(j+1)$ and
$B_M(j)\ge B_{2M}(j)$ [Eq.~\eqref{eq:b-residual}], solving
Eq.~\eqref{eq:b-bridge} for $\Pt^{\opp}(m,\beta_j)$ and dividing the $m=2M$
identity by the $m=M$ identity for even $M$ yields
\begin{equation}
 \frac{\Pt^{\opp}(2M,\beta_j)}{\Pt^{\opp}(M,\beta_j)}
 =r_j^M\underbrace{\frac{B_{2M}(j+1)}{B_M(j+1)}}_{\displaystyle\ge1-e_M(j+1)}
       \underbrace{\left[\frac{B_M(j)}{B_{2M}(j)}\right]^2}_{\displaystyle\ge1}.
 \label{eq:b-coupling}
\end{equation}
Because $r_j^M\approx1$ [Eq.~\eqref{eq:b-ratios}], purity remains controlled upon doubling length.
Squaring the Gibbs operator advances temperature, as in Eq.~\eqref{eq:squaring}:
\begin{equation}
 \tncolumn{\rho_m^{\opp},\rho_m^{\opp}}\;=\;\Pt^{\opp}(m,\beta_j)\,\tncolumn[item={minimum height=1.6em}]{\rho_m^{\opp}(\beta_{j+1})},
 \label{eq:stack}
\end{equation}
so the squaring bound~\eqref{eq:square} at $m=2M$ drives $1-\Pt^{\opp}$ to $\mathcal O(x^2)$.

\emph{Scale induction.---}
Together, the overlap and purity bounds give an induction in which length and
inverse temperature double at each step at the scales
\begin{equation}
 \beta_j=49\,2^{j-2},\quad n_j=120\,2^{j-2},\quad
 M_j=2n_j\quad(j\ge2),
 \label{eq:b-scales}
\end{equation}
with boundary-defect lengths $k_2=142$ and $k_{j+1}=M_j$. Using the ring purity
defect $\varepsilon_j$ ($\varepsilon_2=.052$, $\varepsilon_{j+1}=6\varepsilon_j^2$)
bounded in Appendix~\ref{app:b-periodic} [Eq.~\eqref{eq:b-tail}] and the
starting estimates at $j=2$ in Appendix~\ref{app:b-iteration}, three explicit steps
reach $j=5$, and induction on $j\ge5$ proves
\begin{equation}
 \begin{gathered}
 \varepsilon_j,\ e_{k_j}(j),\ 1-\Pt^{\opp}(M_j,\beta_j)\le x_j\quad(j\ge5),\\
 x_5=.01,\qquad x_{j+1}=10x_j^2,\qquad
 x_j=10^{-1-2^{j-5}}.
 \end{gathered}
 \label{eq:b-x}
\end{equation}
For $j\ge3$, $k_j=M_j/2$, leaving $M_j/2$ sites to suppress remainders.
One inductive step $j\to j+1$ is three moves. In the diagram below, read from
top to bottom, each box is one bound: grey for the boundary defect $e_k(j)$,
blue for the opposite-field thermal purity $\Pt^{\opp}$. Solid boxes are given
at scale $j$ or proved at scale $j+1$, the dashed box is intermediate, and the
grey arrow entering a move is the second bound it uses:
\begin{proofdiagram}
\begin{prooflanes}
 \lanebound{b}{e0}{0}{0}{$e_{M_j/2}(j)\le x$}
 \lanebound{t}{t0}{1}{0}{$\Pt^{\opp}(M_j,\beta_j)\ge1-x$}
 \lanebound{b}{e1}{0}{1}{$e_{M_j}(j{+}1)\le10x^2$}
 \lanebound{to}{t1}{1}{2}{$\Pt^{\opp}(2M_j,\beta_j)\ge1-2x-16x^2$}
 \lanebound{t}{t2}{1}{3}{$\Pt^{\opp}(2M_j,\beta_{j+1})\ge1-10x^2$}
 \lanemove{e0}{e1}{1}{\eqref{eq:b-new-e}} \laneuse{t0}{1}
 \lanemove{t0}{t1}{2}{\eqref{eq:b-coupling}} \laneuse{e1}{2}
 \lanemove[mpoblue]{t1}{t2}{3}{square}
 \lanetag[above]{e0}{given} \lanetag[above]{t0}{given}
 \lanetag[below]{e1}{proved} \lanetag{t2}{proved}
\end{prooflanes}
\end{proofdiagram}
Move~1 is the overlap recurrence~\eqref{eq:b-new-e}--\eqref{eq:b-new-ratio},
move~2 the length doubling~\eqref{eq:b-coupling},
and move~3 the squaring~\eqref{eq:stack}. Since $M_{j+1}=2M_j$, the proved
pair is the given pair of scale $j+1$, so the step repeats. The defect falls doubly exponentially,
$x_5=10^{-2}$\,\sparkline{-17}{-2}{-2,-3,-5,-9,-17}\,$x_9=10^{-17}$ (log scale); since
$\beta_j=392\,2^{j-5}$, it decays exponentially in inverse temperature.

\emph{Equal boundary fields and odd lengths.---}
The opposite-field estimates extend to the equal-field Hamiltonian~\eqref{eq:b-model} via the spin-rotation symmetry $U_P$.
Next to the leading projector, the rotation operator acts trivially:
\begin{equation}
 \tnchain[last=sym]{\Pi_j,U_P}\;=\;\tnchain{\Pi_j}.
 \label{eq:b-twist}
\end{equation}
Thus $Q_j(n):=Y_n^{\same}(\beta_j)/\lambda_j^n$ shares the leading value $c_j$.
For any even $M\le n$, commuting $U_P$ through the right-hand tensors ($[U_P,R_j]=0$) isolates the remainder:
\begin{align}
 Q_j(n)-c_j&=\tnchain[capl=\bar v_j,capr=v_j]{R_j^{M/2},R_j^{n-M}U_P,R_j^{M/2}}\notag\\
 &=\tnchain[capl={\bar u},capr={u}]{R_j^{n-M}U_P},\qquad
 u=R_j^{M/2}v_j.
 \label{eq:b-signed}
\end{align}
Because $\|R_j^{n-M}U_P\|\le1$,
\begin{equation}
 |Q_j(n)-c_j|\le\|u\|^2=D_M(j),
 \label{eq:b-odd-bound}
\end{equation}
bounding the odd-length remainder by $D_M(j)$ without requiring positivity of $R_j^n$.

For $M_j\le n\le2M_j$ and $M=M_j=k_{j+1}$, Eq.~\eqref{eq:b-odd-bound} gives
$Q\le c+D_M=B_M$ and $Q\ge c-D_M=(1-2e_M)B_M$; with $e_M(j+1)\le x_{j+1}$
from Eq.~\eqref{eq:b-x}, this yields
\begin{equation}
 Q_j(n)\le B_M(j),\quad
 Q_{j+1}(n)\ge(1-2x_{j+1})B_M(j+1)>0.
 \label{eq:b-odd-new}
\end{equation}
Since $\Pt^{\same}(n,\beta_j)=Q_{j+1}(n)\lambda_{j+1}^n/[Q_j(n)\lambda_j^n]^2$,
inserting these bounds and Eq.~\eqref{eq:b-bridge} at $m=M$, and using
$r_j^{n-M}\ge r_j^M$ [Eq.~\eqref{eq:b-ratios}] with
$\varepsilon_j\le x_j$, $\varepsilon_{j+1}\le6x_j^2$, $x_{j+1}=10x_j^2$, gives
the equal-field purity:
\begin{align}
 \Pt^{\same}(n,\beta_j)
 &=\frac{Y_n^{\same}(\beta_{j+1})}{Y_n^{\same}(\beta_j)^2}
 \ge(1-2x_{j+1})\Pt^{\opp}(M,\beta_j)r_j^{n-M}\notag\\
 &\ge(1-2x_{j+1})(1-x_j)(1-\varepsilon_j)(1-\varepsilon_{j+1})\notag\\
 &\ge1-2x_j-26x_j^2\ge1-4x_j.
 \label{eq:b-final-purity}
\end{align}

\begin{proof}[Proof of Theorem~\ref{thm:b-gap}]
The intervals $[M_j,2M_j]$ cover every integer $n\ge M_5=1920$, and
Eq.~\eqref{eq:b-final-purity} holds for all such $n$ (Theorem~\ref{thm:b-gap}
states odd $n=2L+1\ge1921$ to match the site-centered reflection setting of
Tasaki's index theorem in Ref.~\cite{boundary}). For any such $n$, choose
$j\ge5$ with $M_j\le n\le2M_j$.
The largest Gibbs weight satisfies $w_0\ge\Pt^{\same}(n,\beta_j)\ge1-4x_j>1/2$:
the spectrum looks like \specbars{.97,.02,.01}, so the ground state is unique.
Equation~\eqref{eq:b-target} gives
\begin{equation}
 e^{-\beta_j\Delta_n}\le\frac{4x_j}{1-4x_j}
 \le\frac{5}{12}\,10^{-2^{j-5}}<10^{-2^{j-5}}.
 \label{eq:b-gap-final}
\end{equation}
Taking logarithms and using $\beta_j=392\,2^{j-5}$ yields $\Delta_n>(\log10)/392$.
\end{proof}

\section{Transfer bounds at the open-chain scales}\label{app:b-periodic}
For the scales in Eq.~\eqref{eq:b-scales}, Appendix~\ref{app:open} uses
\begin{equation}
 \begin{gathered}
 \|R_j\|^{n_j}\le d(\varepsilon_j),\qquad d(\varepsilon)=\frac{\varepsilon}{2-3\varepsilon},\\
 \varepsilon_2=.052,\qquad \varepsilon_{j+1}=6\varepsilon_j^2,
 \end{gathered}
 \label{eq:b-tail}
\end{equation}
and
\begin{equation}
 0<r_j\le1,\quad r_j^{M_j}\ge(1-\varepsilon_j)(1-\varepsilon_{j+1}),\quad
 U_P\Pi_j=\Pi_j.
 \label{eq:b-ratios}
\end{equation}
These follow from the first row of Eq.~\eqref{eq:b-seeds} by the doubling
estimate of the main text, with a larger starting defect: if
$\Ps(n,\beta),\Pt(2n,\beta)\ge1-\varepsilon$, move~1 of Eq.~\eqref{eq:update}
gives $\Ps(n,2\beta)\ge(1-\varepsilon)^3$, and squaring with Eq.~\eqref{eq:square} bounds $1-\Ps(2n,2\beta)$ by
\begin{equation}
 f(1-(1-\varepsilon)^3)=\frac{\varepsilon^2}{2}
 \left[\sum_{s=1}^3(1-\varepsilon)^{-s}\right]^2\le6\varepsilon^2
 \label{eq:b-periodic-six}
\end{equation}
for $0\le\varepsilon\le.052$, in place of $5\varepsilon^2$. Since $6\varepsilon^2\le\varepsilon$,
moves~3 and 4 give the same bound for $1-\Pt(4n,2\beta)$.
Induction therefore proves
$\Ps(n_j,\beta_j),\Pt(M_j,\beta_j)\ge1-\varepsilon_j$.

To turn the purity bound $\Ps(n_j,\beta_j)\ge1-\varepsilon_j$ into the remainder
bound $\|R_j\|^{n_j}\le d(\varepsilon_j)$ in Eq.~\eqref{eq:b-tail}, consider a
probability list with purity at least $1-\varepsilon>1/2$ and largest weight
$p_0$; put $z=1-p_0$. The purity is at most
$p_0^2+z^2=1-2z(1-z)$, so $2z(1-z)\le\varepsilon$. It is also at most $p_0$,
so $p_0>1/2$: the largest weight is unique and $z<1/2$. Then
$z\le2z(1-z)\le\varepsilon$, hence $z\le\varepsilon/[2(1-\varepsilon)]$ and $z/(1-z)\le d(\varepsilon)$.
Applied to the eigenvalues $p_i=|\lambda_i(\beta_j)|^{n_j}/Z_{n_j}(\beta_j)$ of $\sigma_{n_j}(\beta_j)$,
the largest-modulus eigenvalue is simple and
\begin{equation}
 \sum_{i\ge1}\left|\frac{\lambda_i(\beta_j)}{\lambda_j}\right|^{n_j}
 \le d(\varepsilon_j),
 \label{eq:b-full-tail}
\end{equation}
which implies Eq.~\eqref{eq:b-tail}.
Positivity of $Z_N(\beta_j)$ for all large odd $N$ ensures that the leading eigenvalue is positive, since a negative leading eigenvalue would make $\trv X_{\beta_j}^N$ negative.
The residual sum tends to zero since $X_{\beta_j}$ is Hilbert--Schmidt
and its residual spectral radius is strictly smaller than $\lambda_j$.

It remains to show $U_P\Pi_j=\Pi_j$ in Eq.~\eqref{eq:b-ratios}. The real
kernel admits real orthogonal virtual actions of the three $\pi$ rotations and
the cyclic permutation of the coordinate axes, all commuting with
$X_{\beta_j}$. On its real one-dimensional leading eigenspace these actions are
signs. The order-three cyclic rotation has sign $+1$ and conjugates the three
$\pi$ rotations into one another, so their signs agree; since the product of
two is the third, their common sign is $+1$, giving $U_P\Pi_j=\Pi_j$.

Finally, to bound $r_j=\lambda_{j+1}/\lambda_j^2$ in Eq.~\eqref{eq:b-ratios},
note that $Z_N(\beta_{j+1})\le Z_N(\beta_j)^2$ since the thermal purity
$\Pt(N,\beta_j)$ is at most one. Taking $N$th roots along even $N\to\infty$ gives
$\lambda_{j+1}\le\lambda_j^2$. For the lower bound at the required
finite length, the largest eigenvalue of $\sigma_{M_j}(\beta_{j+1})$
is at least $1-\varepsilon_{j+1}$, since $M_j=n_{j+1}$. Hence
\begin{align*}
 \lambda_{j+1}^{M_j}
 &\ge(1-\varepsilon_{j+1})Z_{M_j}(\beta_{j+1})\\
 &\ge(1-\varepsilon_{j+1})(1-\varepsilon_j)Z_{M_j}(\beta_j)^2\\
 &\ge(1-\varepsilon_{j+1})(1-\varepsilon_j)\lambda_j^{2M_j},
\end{align*}
which is the remaining bound in Eq.~\eqref{eq:b-ratios}.


\section{Starting estimates for the open chain}\label{app:b-iteration}
The induction of Appendix~\ref{app:open} [Eq.~\eqref{eq:b-x}] starts at $j=5$,
where all three defects are below $x_5=.01$, whereas the certified inputs are
given at $j=2$. This appendix bridges $j=2$ to $j=5$ in four steps: it states
the inputs, turns the three moves of Appendix~\ref{app:open} into numerical
update rules, runs those rules for three steps, and verifies the induction step
for $j\ge5$.

\emph{Inputs.---}
The certified starting estimates combine periodic variational estimates with
small-cluster boundary certificates from Ref.~\cite{boundary}:
\begin{equation}
 \begin{gathered}
 \Ps(120,49)\ge .948,\qquad \Pt(240,49)\ge .948,\\
 e_{142}(2)\le .14,\qquad \Pt^{\opp}(240,\beta_{2})\ge .79.
 \end{gathered}
 \label{eq:b-seeds}
\end{equation}
The first row gives the ring purity defect $\varepsilon_2=.052$ at
$(n_2,\beta_2)=(120,49)$ used in Appendix~\ref{app:b-periodic}; the second row
gives the open-chain boundary defect at $k_2=142$ and opposite-field thermal
purity at $M_2=240$.

\emph{Update rules.---}
Write $E_j$ for an upper bound on $e_{k_j}(j)$ and $\eta_j$ for an upper bound
on $1-\Pt^{\opp}(M_j,\beta_j)$. Move~1 controls the overlap: combining
Eqs.~\eqref{eq:b-new-e} and \eqref{eq:b-new-ratio} gives
\begin{equation}
 \begin{gathered}
 e_m(j+1)\le F(q,e_k(j),\Pt^{\opp}(m,\beta_j)),\\
 F(q,e,t)=\frac{q}{1-q}\left[\frac1{(1-e)^2t}-1\right],
 \end{gathered}
 \label{eq:b-F}
\end{equation}
which increases with $q,e$ and decreases with $t$. The suppression factor $q$
must bound $\|R_{j+1}\|^{M_j-k_j}$; since Eq.~\eqref{eq:b-tail} bounds
$\|R_{j+1}\|^{n_{j+1}}$ and $n_{j+1}=M_j$, this fixes the exponent
$(M_j-k_j)/M_j$ in Eq.~\eqref{eq:b-q}. Moves~2 and 3 control the purity via
Eqs.~\eqref{eq:b-coupling} and \eqref{eq:square}. With $k_2=142$ and
$k_{j+1}=M_j$, the updates are
\begin{align}
 q_j&=d(\varepsilon_{j+1})^{(M_j-k_j)/M_j},\label{eq:b-q}\\
 E_{j+1}&\ge F(q_j,E_j,1-\eta_j),\label{eq:b-E-update}\\
 \eta_{j+1}&\ge f\bigl(1-(1-\eta_j)(1-\varepsilon_j)
                   (1-\varepsilon_{j+1})(1-E_{j+1})\bigr),
 \label{eq:b-eta-update}
\end{align}
for any $E_{j+1},\eta_{j+1}\in[0,1)$ satisfying these inequalities.

\emph{Three explicit steps.---}
Starting from $\varepsilon_2=.052$, $E_2=.14$, and $\eta_2=1-.79=.21$, three
applications of Eqs.~\eqref{eq:b-E-update}--\eqref{eq:b-eta-update} give the
rational upper bounds:
\begin{center}
\begin{tabular}{@{}rrrr@{}}
\toprule
$j$ & $E_j$ & $\eta_j$ & $q_j$ (upper bound)\\
\midrule
2 & .14 & .21 & .142\\
3 & .12 & .155 & .0282\\
4 & .016 & .026 & .0028\\
5 & .001 & .001 & ---\\
\bottomrule
\end{tabular}
\end{center}
Both defects collapse along the iteration, $E_j$:
\specbars[black!45]{1,.86,.11,.007} and $\eta_j$: \specbars{1,.74,.12,.005}
($j=2,\dots,5$, each relative to $j=2$).
Rational arithmetic confirms these step bounds (e.g., at $j=2$ the exponent in
Eq.~\eqref{eq:b-q} is $(240-142)/240=49/120$, so $q_2\le.142$ amounts to
$d(\varepsilon_3)^{49}<.142^{120}$, with subsequent steps verified by
squaring), and $\varepsilon_5<.000015$, establishing the base case $j=5$ of
Eq.~\eqref{eq:b-x}.

\emph{The induction step.---}
Assuming Eq.~\eqref{eq:b-x} at scale $j\ge5$ with $x=x_j\le.01$, the relation
$k_j=M_j/2$ gives
\begin{equation}
 \begin{gathered}
 q_j\le\sqrt{\frac{6x^2}{2-18x^2}}\le2x,\\
 F(2x,x,1-x)\le\frac{6x^2}{(1-2x)(1-x)^3}<7x^2.
 \end{gathered}
 \label{eq:b-x-e}
\end{equation}
Setting $E_{j+1}=10x^2$ satisfies Eq.~\eqref{eq:b-E-update}.
The four-factor product in Eq.~\eqref{eq:b-eta-update} is at least $1-2x-16x^2$, and $f(2x+16x^2)<3x^2$, so $\eta_{j+1}=10x^2$ satisfies Eq.~\eqref{eq:b-eta-update} with $\varepsilon_{j+1}\le6x^2$, completing the induction for Eq.~\eqref{eq:b-x}.


\end{document}
